\section*{Sets and Indices}

Let
\[
\mathcal{P}=\left\{(a,b)\in \mathbb{Z}_{+}^{2} \;:\; 3a+4b \le L,\; a+b \ge 1\right\}
\]
be the set of all feasible cutting patterns generated by the code, where each pattern specifies how many 3-meter and 4-meter pieces are cut from one raw bar.

Index patterns by \(p \in \mathcal{P}\). For a pattern \(p\), let
\[
a_p = \text{number of 3-meter pieces in pattern } p,\qquad
b_p = \text{number of 4-meter pieces in pattern } p.
\]

\section*{Parameters}

From the CSV data:

\begin{itemize}
    \item \(q_3\): required number of 3-meter bars (code: \texttt{q3}); here \(q_3=90\).
    \item \(q_4\): required number of 4-meter bars (code: \texttt{q4}); here \(q_4=60\).
    \item \(L\): raw bar length in meters (code: \texttt{L}); here \(L=10\).
\end{itemize}

Fixed piece lengths used in the implementation:

\begin{itemize}
    \item \(\ell_1 = 3\): length of type-1 bar (code: \texttt{len1}).
    \item \(\ell_2 = 4\): length of type-2 bar (code: \texttt{len2}).
\end{itemize}

\section*{Decision Variables}

\begin{itemize}
    \item \(x_p \in \mathbb{Z}_{+}\): number of raw bars cut according to pattern \(p\) (code: \texttt{vars[i]}).
\end{itemize}

\section*{Objective}

The implemented objective minimizes total raw length used minus the total required length:
\[
\min \sum_{p \in \mathcal{P}} L\,x_p - (\ell_1 q_3 + \ell_2 q_4).
\]

Equivalently, since the second term is constant, this is the same implemented optimization as minimizing the number of raw bars used:
\[
\min \sum_{p \in \mathcal{P}} x_p.
\]

\section*{Constraints}

Demand satisfaction is enforced with inequalities, so overproduction is allowed:

\[
\sum_{p \in \mathcal{P}} a_p x_p \ge q_3
\]

\[
\sum_{p \in \mathcal{P}} b_p x_p \ge q_4
\]

Integrality and nonnegativity:
\[
x_p \in \mathbb{Z}_{+} \qquad \forall p \in \mathcal{P}.
\]

\section*{Notes on Implementation}

The code explicitly enumerates all feasible single-bar cutting patterns
\[
(a,b)\in \mathbb{Z}_{+}^{2} \text{ such that } 3a+4b\le L,\; (a,b)\neq(0,0),
\]
and then solves the resulting integer linear program.

For this instance with \(L=10\), the enumerated patterns are:
\[
(0,1),\ (0,2),\ (1,0),\ (1,1),\ (2,0),\ (2,1),\ (3,0).
\]

The objective is described in the code comments as ``total waste,'' computed as total raw-bar length used minus the required total piece length. Because the required total piece length is a constant, any overproduction is not charged by piece length directly; it affects the objective only through the additional raw bars required. The demand constraints are \(\ge\) constraints exactly as implemented, not equalities. The model is solved as an integer program using PuLP with the Gurobi command-line solver backend.
